\chapter{Lattice and Band Geometry}
\label{ch:lattice_band}

\section{Introduction and Discrete Setting}

In $1+1$-dimensional field theory, continuum descriptions can produce ultraviolet and infrared divergences, as well as infinite mode sums that require separate regularization. Here we begin with a finite periodic chain and its exact Fourier-dual momentum band. The resulting operators and sums are finite, and the algebraic constraints imposed by periodicity can be handled before any continuum limit is considered.

Throughout this chapter, let $L \in \mathbb{N}$ denote the number of spatial lattice sites, assuming $L \ge 1$. While general difference calculus and Fourier transforms can be formulated for arbitrary positive integers $L$, the construction of half-filled Dirac seas and fermionic Fock spaces from Chapter~\ref{ch:fermions} onward strictly requires an even lattice length:
\begin{equation}
L = 2h, \quad h \in \mathbb{N}, \quad h \ge 1.
\end{equation}
The integer $h$ represents the half-filling particle number or the positive Nyquist cut-off.

\section{The Spatial Lattice}

The spatial geometry of the physical system is that of a discrete one-dimensional circular chain of $L$ sites with periodic boundary conditions.

\begin{definition}[Spatial Lattice]
The spatial lattice $\Lambda$ is defined as the cyclic quotient group of integers modulo $L$:
\begin{equation}
\Lambda \coloneqq \mathbb{Z}/L\mathbb{Z}.
\end{equation}
Elements $x, y \in \Lambda$ represent discrete spatial positions. Periodic boundary conditions are an intrinsic algebraic property of the quotient ring: translation by $L$ sites corresponds to the identity operation in $\Lambda$.
\end{definition}

Because $L \ge 1$, the spatial lattice $\Lambda$ is a finite commutative ring of cardinality $\#\Lambda = L$, endowed with canonical addition and multiplication inherited from $\mathbb{Z}$.

\begin{remark}[Finite-Lattice Remark]
The quotient description makes periodicity exact: adding $L$ to an integer representative leaves its site unchanged. Since $L$ is positive, this finite ring has exactly $L$ sites, and all position labels and translations are interpreted within it.
\end{remark}

\paragraph{Formal Correctness vs.\ Physical Implications: Discrete Chains vs.\ Continuum Circles.}
In condensed matter physics, a circular chain is often treated in the continuum limit by taking the lattice spacing $a \to 0$ and the number of sites $L \to \infty$ while holding the macroscopic circumference $\mathcal{L} = L a$ fixed. In that limit, spatial positions $x a$ tend to continuous coordinates on the circle $\mathbb{R}/\mathcal{L}\mathbb{Z}$, and fields become continuous functions or operator-valued distributions. 
In our exact algebraic framework, by contrast, the lattice spacing $a$ is set to $1$ without loss of generality, and all positions $x \in \Lambda$ are exact residue classes. This eliminates distribution-theoretic smearing and guarantees that every spatial operator algebra is strictly finite-dimensional.

\section{The Momentum Band and Nyquist Asymmetry}

In continuous quantum field theory, the dual momentum space is the real line $\mathbb{R}$ or the continuous circle $\mathbb{R}/(2\pi/a)\mathbb{Z}$. On a discrete spatial lattice of $L$ sites, the dual momentum band (the first Brillouin zone) must consist of exactly $L$ discrete momentum modes.

To center the physical momentum around zero while maintaining exact finite cardinality, we define the momentum band as an integer interval centered near the origin.

\begin{definition}[Momentum Band]
The momentum band $\Lambda^*$ is the subset of signed integers satisfying the two-sided inequalities:
\begin{equation}
\Lambda^* \coloneqq \left\{ k \in \mathbb{Z} \;\middle|\; -L < 2k \le L \right\}.
\end{equation}
Momenta are denoted by variables $k, p, q \in \Lambda^*$. The physical wavenumber corresponding to mode index $k$ is given by the scaled wavenumber:
\begin{equation}
p_{\mathrm{phys}} \coloneqq \frac{2\pi k}{L}.
\end{equation}
\end{definition}

The choice of inequalities $-L < 2k \le L$ is subtle and fundamental:
\begin{itemize}
    \item When $L = 2h$ is even, the condition $-2h < 2k \le 2h$ simplifies to $-h < k \le h$. Consequently, the momentum band consists explicitly of the $2h$ consecutive integers:
    \begin{equation}
    \Lambda^* = \{-h+1, -h+2, \dots, -1, 0, 1, \dots, h-1, h\}.
    \end{equation}
    \item The interval is intrinsically asymmetric: it includes the positive Nyquist mode $k = +h = +L/2$, but excludes the negative counterpart $k = -h = -L/2$.
    \item When $L$ is odd, say $L = 2m+1$, the condition becomes $-2m-1 < 2k \le 2m+1$, which yields the symmetric set $\{-m, -m+1, \dots, m-1, m\}$ of cardinality $2m+1$.
\end{itemize}

\paragraph{Formal Correctness vs.\ Physical Implications: Brillouin Zone Boundaries.}
In standard solid-state physics textbooks (such as Ashcroft and Mermin or Kittel), the first Brillouin zone is customarily drawn as the symmetric interval $[-\pi/a, \pi/a]$, or $[-h, h]$ in mode units. However, on an even lattice of $L = 2h$ sites, the modes $k = -h$ and $k = +h$ represent the \emph{same} physical wave on the lattice, since $e^{i (-h) (2\pi x / 2h)} = e^{-i \pi x} = (-1)^x = e^{i \pi x} = e^{i (+h) (2\pi x / 2h)}$.
Including both $-h$ and $+h$ yields $2h+1$ modes, which violates the fundamental dimension theorem $\dim(\ell^2(\Lambda)) = \dim(\ell^2(\Lambda^*)) = L$. Physics literature often circumvents this by assigning fractional weights $1/2$ to the boundary modes or passing implicitly to the thermodynamic limit. 
The discrete Fourier transform pairs spaces of the same finite dimension $L$. Retaining the positive Nyquist mode $+h$ while excluding $-h$ gives exactly $\#\Lambda^* = L$ distinct periodic modes.

\begin{remark}[Finite-Lattice Remark]
The band is an explicitly bounded set of integer modes, with exactly one representative of each residue class modulo $L$. Its asymmetric endpoints record the single positive Nyquist mode on an even lattice.
\end{remark}

\section{Quotient Projections and the Centered Section}

Because $\Lambda = \mathbb{Z}/L\mathbb{Z}$ is a quotient group, there is a natural projection from any integer to its equivalence class in the lattice.

\begin{definition}[Canonical Projection]
The canonical projection $\pi : \mathbb{Z} \to \Lambda$ maps an integer $k \in \mathbb{Z}$ to its periodic residue class:
\begin{equation}
\pi(k) \coloneqq k \bmod L.
\end{equation}
By restriction, $\pi$ maps any momentum mode $k \in \Lambda^*$ to an element $\pi(k) \in \Lambda$.
\end{definition}

To construct the inverse mapping, we seek a section of $\pi$ whose range lies strictly within the momentum band $\Lambda^*$. Every element $x \in \Lambda = \mathbb{Z}/L\mathbb{Z}$ possesses a unique canonical representative $x_{\mathrm{val}} \in \{0, 1, \dots, L-1\}$.

\begin{definition}[Centered Lift Section]
The centered lift map $r : \Lambda \to \Lambda^*$ is defined piecewise on the unsigned representative $x_{\mathrm{val}} \in \{0, 1, \dots, L-1\}$ by:
\begin{equation}
r(x) \coloneqq \begin{cases}
x_{\mathrm{val}}, & \text{if } 2 x_{\mathrm{val}} \le L, \\
x_{\mathrm{val}} - L, & \text{if } 2 x_{\mathrm{val}} > L.
\end{cases}
\end{equation}
\end{definition}

\paragraph{Physical Intuition for the Centered Lift.}
The map $r$ takes a non-negative residue $x_{\mathrm{val}} \in [0, L-1]$ and centers it around zero: values in the upper half of the range $[h+1, 2h-1]$ are translated downward by $L$ to become negative integers in $[-h+1, -1]$, corresponding to negative-frequency or left-moving momentum modes in physical terminology.

\begin{lemma}[Well-Definedness and Inversion]
\label{lem:section_identities}
For any positive integer $L \ge 1$:
\begin{enumerate}
    \item For every $x \in \Lambda$, the value $r(x)$ belongs to the momentum band $\Lambda^*$, and:
    \begin{equation}
    \pi(r(x)) = x.
    \end{equation}
    \item For every $k \in \Lambda^*$, the centered lift recovers the original mode:
    \begin{equation}
    r(\pi(k)) = k.
    \end{equation}
\end{enumerate}
\end{lemma}

\begin{proof}
We verify both statements directly from the definitions:
\begin{enumerate}
    \item Let $x \in \Lambda$ and let $a = x_{\mathrm{val}} \in \{0, 1, \dots, L-1\}$.
    \begin{itemize}
        \item If $2a \le L$, then $r(x) = a$. Since $a \ge 0$, we have $2a \ge 0 > -L$. Since $2a \le L$, $a$ satisfies $-L < 2a \le L$, so $r(x) \in \Lambda^*$. Furthermore, $\pi(r(x)) = \pi(a) = x$.
        \item If $2a > L$, then $r(x) = a - L$. Since $a \le L - 1$, we have $a - L \le -1$, so $2(a - L) \le -2 < 0 \le L$. Conversely, $2a > L$ implies $2(a - L) > -L$. Thus $-L < 2(a - L) \le L$, establishing that $r(x) \in \Lambda^*$. Moreover, $\pi(a - L) = \pi(a) - \pi(L) = x - 0 = x$.
    \end{itemize}
    In both branches, $r(x) \in \Lambda^*$ and $\pi(r(x)) = x$.
    
    \item Let $k\in\Lambda^*$. By part (1), $k$ and $r(\pi(k))$ have the same image under $\pi$, so their difference is an integer multiple of $L$. Since both lie in the band $(-L/2,L/2]$, their difference has absolute value strictly less than $L$. The only such multiple is zero, hence $r(\pi(k))=k$.
\end{enumerate}
\end{proof}

\begin{corollary}[Canonical Bijection]
The restriction of the canonical projection $\pi|_{\Lambda^*} : \Lambda^* \to \Lambda$ is a bijection, whose two-sided inverse is the centered lift section $r : \Lambda \to \Lambda^*$. In particular, the cardinality of the momentum band is strictly:
\begin{equation}
\#\Lambda^* = \#\Lambda = L.
\end{equation}
\end{corollary}

\begin{remark}[Finite-Lattice Remark]
For injectivity, take any $k_1,k_2\in\Lambda^*$ with $\pi(k_1)=\pi(k_2)$. Their difference is an integer multiple of $L$. Both lie in $(-L/2,L/2]$, so the diameter of the band is strictly less than $L$ and $|k_1-k_2|<L$. The difference must therefore be zero. Thus $k_1=k_2$, and the section identities give the asserted inversion.
\end{remark}

\section{Transported Band Arithmetic and Umklapp Wrapping}

While $\Lambda = \mathbb{Z}/L\mathbb{Z}$ is an abelian group under addition, the momentum band $\Lambda^*$ is an integer subset. The usual sum of two momenta $k, p \in \Lambda^*$ may lie outside the interval $(-L/2, L/2]$. 

To endow $\Lambda^*$ with a group structure isomorphic to $\Lambda$, we transport the addition and negation of the cyclic group through the bijection $(\pi, r)$.

\begin{definition}[Transported Band Arithmetic]
For any momenta $k, p \in \Lambda^*$, we define band addition $\oplus$ and band negation $\ominus$ by:
\begin{align}
k \oplus p &\coloneqq r(\pi(k) + \pi(p)), \\
\ominus k &\coloneqq r(-\pi(k)).
\end{align}
\end{definition}

\begin{remark}[Notation in later chapters]
The symbols $\oplus$ and $\ominus$ emphasize that these operations are transported from the quotient group to the chosen momentum band. In subsequent chapters, standard $+$ and $-$ symbols will denote these band operations when their domain is clear, reducing notational overhead.
\end{remark}

Because $(\pi, r)$ is a bijection between $\Lambda^*$ and the abelian group $\Lambda$, $(\Lambda^*, \oplus, \ominus, 0)$ forms a finite abelian group isomorphic to $\mathbb{Z}/L\mathbb{Z}$, with neutral element $0 = r(\pi(0))$.

\paragraph{Physical Intuition: Umklapp Scattering on a Lattice.}
In solid-state physics, when two quasi-particles with crystal momenta $k$ and $p$ scatter, their total momentum $k+p$ can exceed the boundary of the first Brillouin zone. The crystal lattice absorbs a reciprocal lattice vector $G = \pm L$, scattering the particles back into the first Brillouin zone with crystal momentum $k \oplus p = k + p \mp L$. This physical phenomenon is known as an \emph{Umklapp process}. Theorem~\ref{thm:wrap_lemma} makes this exact algebraically: the integer difference between arithmetic addition and band addition is always $w L$ with $w \in \{-1, 0, 1\}$.

\begin{theorem}[Umklapp Wrapping Lemma]
\label{thm:wrap_lemma}
For any $k, p \in \Lambda^*$, there exists a unique integer $w \in \{-1, 0, 1\}$ such that:
\begin{equation}
k + p = (k \oplus p) + w L.
\end{equation}
The integer $w$ is the Umklapp multiplier and is explicitly determined by:
\begin{equation}
w = \begin{cases}
+1, & \text{if } 2(k+p) > L, \\
0, & \text{if } -L < 2(k+p) \le L, \\
-1, & \text{if } 2(k+p) \le -L.
\end{cases}
\end{equation}
\end{theorem}

\begin{proof}
Let $s=k+p$. Since $k,p\in\Lambda^*$, their band bounds give $-L<s\le L$. By definition, $k\oplus p$ and $s$ have the same image under $\pi$, so $s-(k\oplus p)=wL$ for some integer $w$. The bounds $-L<s\le L$ and $-L/2<k\oplus p\le L/2$ imply $-3L/2<s-(k\oplus p)<3L/2$, hence $-3/2<w<3/2$. Therefore $w\in\{-1,0,1\}$.
By the defining property of the centered lift, $k\oplus p=r(\pi(s))$ is the unique element of $\Lambda^*$ with quotient image $\pi(s)$. Thus the integer $w$ above is uniquely determined by $s-(k\oplus p)=wL$. Its value must be one of $-1,0,1$, as established by the bounds.
Uniqueness of $w$ follows since $(w_1 - w_2) L = 0$ implies $w_1 = w_2$ when $L > 0$.
\end{proof}

\section{The Exceptional Nyquist Mode}

We now examine the behavior of band negation $\ominus k$ when the lattice length is even, $L = 2h$.

\begin{theorem}[Nyquist Mode Negation]
\label{thm:nyquist_neg}
Let $L = 2h$ with $h \ge 1$.
\begin{enumerate}
    \item The positive Nyquist mode $k = h = L/2$ is self-inverse under band negation:
    \begin{equation}
    \ominus \left(\frac{L}{2}\right) = \frac{L}{2}.
    \end{equation}
    \item For any mode $k \in \Lambda^*$ distinct from the Nyquist mode ($k \neq L/2$), band negation coincides with ordinary integer arithmetic negation:
    \begin{equation}
    \ominus k = -k.
    \end{equation}
\end{enumerate}
\end{theorem}

\begin{proof}
Let $L = 2h$.
\begin{enumerate}
    \item For the mode $k = L/2 = h$, ordinary negation yields $-k = -h = -L/2$. In the quotient ring $\Lambda = \mathbb{Z}/L\mathbb{Z}$:
    \begin{equation}
    -\pi(L/2) = \pi(-L/2) = \pi(-L/2 + L) = \pi(L/2).
    \end{equation}
    Applying the section $r$, we have:
    \begin{equation}
    \ominus (L/2) = r(-\pi(L/2)) = r(\pi(L/2)) = L/2,
    \end{equation}
    by Lemma~\ref{lem:section_identities}(2). Hence $\ominus(L/2) = L/2$.

    \item Suppose $k \in \Lambda^*$ with $k \neq L/2$. Since $k \in \{-h+1, \dots, h\}$ and $k \neq h$, we must have:
    \begin{equation}
    -h+1 \le k \le h-1 \implies -(h-1) \le -k \le h-1.
    \end{equation}
    Multiplying by $2$:
    \begin{equation}
    -2h < -2(h-1) \le 2(-k) \le 2(h-1) < 2h = L.
    \end{equation}
    Thus $-k$ satisfies the strict bounds $-L < 2(-k) \le L$, establishing that the integer $-k$ is already in $\Lambda^*$. By Lemma~\ref{lem:section_identities}(2), $r(\pi(-k)) = -k$. Therefore:
    \begin{equation}
    \ominus k = r(-\pi(k)) = r(\pi(-k)) = -k.
    \end{equation}
\end{enumerate}
\end{proof}

\begin{remark}[Finite-Lattice Remark]
The distinction between band negation and ordinary integer negation matters at the boundary: the positive Nyquist mode is self-inverse, while ordinary negation would leave the chosen band. This exact accounting is important in Fourier inversion and in transformations involving paired modes. Specifically:
\begin{itemize}
    \item The pair $(k, \ominus k)$ contains two distinct modes for all $k \notin \{0, L/2\}$.
    \item The fixed points of the negation involution $k \mapsto \ominus k$ on $\Lambda^*$ are precisely the two zero-energy/boundary modes: $k = 0$ and $k = L/2$.
\end{itemize}
This exact accounting ensures that no modes are double-counted or omitted when momentum sums are partitioned into pairs in later chapters.
\end{remark}

\section{Chapter Summary and Literature Comparison}

In this opening chapter, we established the discrete spatial and momentum geometry that underlies all subsequent field-theoretic constructions:
\begin{itemize}
    \item The spatial lattice is the finite cyclic group $\Lambda = \mathbb{Z}/L\mathbb{Z}$ ($L = 2h$).
    \item The momentum band is the asymmetric Brillouin zone $\Lambda^* = \{-h+1, \dots, h\}$, retaining $+h$ and excluding $-h$.
    \item The canonical projection $\pi$ and centered lift section $r$ form a two-sided bijection $\Lambda^* \cong \Lambda$.
    \item Group arithmetic on $\Lambda^*$ is transported via $(\pi, r)$, with integer discrepancies governed by the Umklapp wrapping multiplier $w \in \{-1, 0, 1\}$.
    \item The positive Nyquist mode $k = h$ is the unique non-zero self-inverse mode under band negation $\ominus$.
\end{itemize}

\paragraph{Comparison with Standard Continuum Literature.}
In standard condensed matter bosonization texts, such as Giamarchi~\cite{giamarchi2004}, Gogolin, Nersesyan, and Tsvelik~\cite{gogolin1998}, and von Delft and Schoeller~\cite{vondelft1998}, the spatial circle is typically taken as a continuous interval $[0, \mathcal{L}]$ with periodic boundary conditions, leading to an infinite set of discrete momentum modes $k_n = \frac{2\pi n}{\mathcal{L}}$ ($n \in \mathbb{Z}$). In that setting, there is no upper momentum cutoff, Umklapp scattering is treated as an extrinsic perturbation or ignored in low-energy Tomonaga--Luttinger models, and the Fermi sea has infinite depth.

By contrast, our finite-lattice formulation models the physical lattice from the outset. As emphasized by Haldane~\cite{haldane1981}, the lattice acts as an intrinsic ultraviolet regulator. The asymmetric Brillouin zone $\Lambda^*$ is the unique choice that simultaneously achieves:
\begin{enumerate}
    \item Exact finite dimension: $\dim(\ell^2(\Lambda)) = \dim(\ell^2(\Lambda^*)) = L$.
    \item Absence of duplicate physical modes.
    \item Unambiguous arithmetic at the band boundaries.
\end{enumerate}
In later chapters (specifically Chapter~\ref{ch:heisenberg_algebra} and Chapter~\ref{ch:haldane_completeness}), we will demonstrate that when particle excitations are restricted to an energy budget margin $2M + K + |N| \le h$, the Umklapp wrapping $w$ vanishes identically ($w = 0$), and the exact lattice dynamics reproduces the continuum current algebra of Giamarchi~\cite{giamarchi2004} and von Delft--Schoeller~\cite{vondelft1998} with zero error.
